\section{Introduction}

Structural optimisation is a cornerstone of modern computational engineering, enabling the design of lighter, more efficient, and more cost-effective structures. 
Among the most widely studied civil structures in this context is the natural-draught cooling tower, whose hyperboloidal shell geometry presents a rich landscape of continuous, nonlinear optimisation challenges~\cite{mungan1976buckling}. 
Minimising the lateral surface area of such a tower, directly proportional to construction material consumption, while satisfying a strict volumetric capacity constraint is a problem of both theoretical interest and direct industrial relevance.

This report constitutes Task 6 of the Digital Engineering Design Group Project undertaken by Group Chenonceau within the MSc in Computational and Software Techniques in Engineering at Cranfield University. 
The group comprises John Hoarau (Project Lead and Integration), Martin Harambat (BFGS Implementation), Pierre B\'{e}drune (Simulated Annealing), and Cl\'{e}mence-Philom\`{e}ne Hinot (Particle Swarm Optimisation and \LaTeX).

\subsection{Project Context and Progression}

Over the preceding tasks, we as a group progressively built our optimisation competence through a structured sequence of increasingly complex problems. 
In Task 3, two metaheuristic algorithms, Particle Swarm Optimisation (PSO) and Simulated Annealing (SA), were independently developed and validated against
 two canonical benchmarks: the discrete Gear Train Design problem and the multi-objective Two-Bar Plane Truss Design problem. 
Task 4 then introduced a deterministic gradient-based counterpart, the Broyden-Fletcher-Goldfarb-Shanno (BFGS) quasi-Newton method, applied to the same benchmarks to
 enable direct algorithmic comparison. 
Key outcomes included the characterisation of BFGS superlinear convergence properties, the evaluation of noise robustness across all three methods, and the identification 
of the complementary strengths and weaknesses of stochastic versus deterministic approaches.

Task 6 represents the culmination of this progression. For the first time, all three algorithms are deployed together on a single, real-world structural optimisation problem: the minimisation of the lateral surface area of a discretised hyperboloidal cooling tower subject to a fixed volume constraint.

\subsection{Problem Statement}

The cooling tower is modelled as a vertical stack of $m$ conical frustums, yielding a constrained nonlinear programme (NLP) 
in which the decision variables, subsets of the internal radii and section heights, are optimised to minimise the total lateral surface area while preserving a target 
internal volume $V_{\text{target}} = 70{,}320\,\text{m}^3$. The equality constraint is handled via a quadratic penalty method,
 enabling uniform treatment across all three solvers without introducing algorithm-specific constraint logic.

To comprehensively evaluate algorithmic performance, eight study cases of increasing complexity are defined. These range from a fixed-height reference case ($D = 9$ decision variables) 
to a parametric hyperboloidal fit ($D = 3$), encompassing variations in decision variable type, dimensionality, initialisation strategy, constraint topology, and physical scaling.

\subsection{Report Structure}

The remainder of this report is organised as follows. Section 2 presents the mathematical formulation of the cooling tower problem, 
the quadratic penalty methodology, the software architecture, and the rationale behind the eight study cases. 
Section 3 details the three optimisation algorithms and their respective adaptations to this problem. 
Section 4 reports the results across all cases and algorithms. 
Section 5 provides a comparative discussion on convergence behaviour, constraint satisfaction, and computational cost. Finally, 
Section 6 draws conclusions and identifies directions for future work.

\subsection{Statement on the Use of Generative AI}

During the preparation of this work, the authors used ChatGPT (OpenAI) and Claude AI (Anthropic) to assist with code efficiency optimisation and debugging, 
technical writing refinement, and literature review formatting. 
After using these tools, all content was reviewed and edited as needed. 
The authors take full responsibility for the final report's content, analysis, and code functionality.
